\documentclass[12pt,letterpaper]{article}

\usepackage[spanish,es-tabla]{babel}
\usepackage[utf8]{inputenc}
\usepackage[T1]{fontenc}
\usepackage{newtxtext,newtxmath}
\usepackage{geometry}
\usepackage{setspace}
\usepackage{amsmath,amssymb}
\usepackage{booktabs}
\usepackage{array}
\usepackage{enumitem}
\usepackage{float}
\usepackage{longtable}
\usepackage{microtype}
\usepackage{hyperref}
\usepackage{fancyhdr}
\usepackage{caption}

\geometry{margin=2.54cm}
\setstretch{1.15}
\setlength{\parindent}{1.27cm}
\setlength{\parskip}{0pt}
\captionsetup{font=small}

\hypersetup{
	colorlinks=true,
	linkcolor=black,
	urlcolor=blue,
	pdftitle={Optimización multiperíodo del abastecimiento farmacéutico en Droguería Rebarata, Piendamó, Cauca},
	pdfauthor={Emanuel Esteban Restrepo Patarroyo y Ariana Marcela Andrade Bello}
}

\pagestyle{fancy}
\fancyhf{}
\fancyhead[R]{\thepage}
\renewcommand{\headrulewidth}{0pt}
\setlength{\headheight}{14.5pt}

\begin{document}

\begin{titlepage}
	\centering
	\vspace*{1.8cm}
	{\Large \textbf{Optimización multiperíodo del abastecimiento farmacéutico minorista en Piendamó mediante programación lineal entera mixta}}\\[0.6cm]
	{\large Reporte ejecutivo bajo rúbrica IMRAD adaptada}\\[1.4cm]

	Emanuel Esteban Restrepo Patarroyo\\
	Ariana Marcela Andrade Bello\\
	Ingeniería de Sistemas -- UNICOMFACAUCA\\
\texttt{esteban.restrepo@unicomfacauca.edu.co}\\
\texttt{ariana.andrade@unicomfacauca.edu.co}\\[0.8cm]

	Curso: Investigación de Operaciones\\
	Docente: Alejandro Benítez Garzón\\[0.8cm]

	\vfill
	{\large Popayán, Colombia}\\
	{\large 12 de mayo de 2026}
\end{titlepage}

\tableofcontents
\listoftables
\newpage

\section*{Resumen}
\addcontentsline{toc}{section}{Resumen}
Este reporte aborda el problema real de abastecimiento de la Droguería Rebarata (Piendamó, Cauca), donde coexisten faltantes en medicamentos críticos y sobreinventario en referencias de baja rotación, con impacto en caja, costos logísticos y continuidad de servicio. El objetivo fue formular y evaluar un modelo de optimización multiproducto y multiperíodo, basado en programación lineal entera mixta (MILP), para definir cantidades óptimas de pedido por período bajo restricciones de demanda, balance de inventario y capacidad de almacenamiento. La metodología adaptó el enfoque de Salas-Navarro et al. (2025) al contexto local, estructuró parámetros por producto y mes, y validó la formulación mediante cuatro pruebas: consistencia dimensional, casos extremos, convergencia entre herramientas (PuLP vs. implementación comparada) y coherencia operativa. En el escenario base (dos productos, tres períodos), el costo óptimo estimado fue de 1{,}850 COP con cumplimiento total de demanda y sin inventario residual; además, el análisis de sensibilidad indicó precio sombra nulo de capacidad en el caso base y valor marginal positivo cuando se impone inventario de seguridad. Frente al esquema reactivo de referencia, se proyecta una mejora de 11.8\% en costo de planeación y mayor estabilidad de abastecimiento. Se recomienda implementar un piloto trimestral para antibióticos y analgésicos críticos, con revisión mensual de parámetros y escalamiento al portafolio completo condicionado a mantener nivel de servicio tipo I y II iguales o superiores a 0.98.

\section*{Abstract}
\addcontentsline{toc}{section}{Abstract}
This report addresses a real replenishment problem at Droguería Rebarata (Piendamó, Cauca, Colombia), where stockouts of critical medicines coexist with overstocks of low-rotation products, affecting cash flow, logistics costs, and service continuity. The objective was to formulate and evaluate a multi-product, multi-period optimization model based on mixed-integer linear programming (MILP) to determine optimal order quantities by period under demand, inventory-balance, and storage-capacity constraints. The methodology adapted the approach of Salas-Navarro et al. (2025) to the local context, structured product-period parameters, and validated the formulation through four checks: dimensional consistency, extreme-case replication, cross-tool convergence (PuLP vs. compared implementation), and operational coherence. In the base scenario (two products, three periods), the estimated optimal cost was COP 1,850 with full demand satisfaction and zero ending inventory; sensitivity analysis also showed a zero capacity shadow price in the base case and a positive marginal value when safety stock is enforced. Compared with the current reactive policy, an 11.8\% planning-cost improvement and more stable replenishment are projected. A three-month pilot is recommended for critical antibiotics and analgesics, with monthly parameter updates and portfolio-wide scaling conditioned on maintaining Type I and Type II service levels at or above 0.98.

\section*{Palabras clave}
\addcontentsline{toc}{section}{Palabras clave}
Optimización multiperíodo; inventarios farmacéuticos; programación lineal entera mixta; incertidumbre de demanda; Piendamó, Cauca.

\section*{Keywords}
\addcontentsline{toc}{section}{Keywords}
Multi-period optimization; pharmaceutical inventory; mixed-integer linear programming; demand uncertainty; Piendamó, Cauca.

\section{Introducción}
La gestión de inventario farmacéutico en municipios intermedios enfrenta un problema operativo recurrente: decisiones de compra reactivas que generan simultáneamente quiebres de inventario en medicamentos trazadores y acumulación de unidades de baja rotación. En la Droguería Rebarata, este patrón incrementa el costo total y deteriora la continuidad de atención. Como línea base de gestión, se toma una política de reposición no optimizada (pedido por agotamiento), contra la cual se evaluará el desempeño económico del modelo propuesto.

La literatura reporta que los modelos de optimización multiperíodo, en particular MILP con extensión estocástica, son adecuados para balancear costo y disponibilidad en cadenas farmacéuticas bajo restricciones reales de capacidad y demanda (Franco \& Alfonso-Lizarazo, 2020; Zandkarimkhani et al., 2020; Salas-Navarro et al., 2025). En consecuencia, el enfoque adoptado prioriza formulación determinística como base y análisis de incertidumbre como extensión gerencial.

El objetivo general de este reporte es \textbf{formular, verificar y evaluar} un modelo MILP multiproducto y multiperíodo para mejorar el abastecimiento de la Droguería Rebarata en Piendamó. El reporte se organiza así: la Sección 2 presenta la revisión de literatura; la Sección 3 describe caso, formulación, verificación y alternativas; la Sección 4 reporta resultados y sensibilidad; la Sección 5 desarrolla análisis, discusión y conclusión con recomendación gerencial.

\section{Revisión de Literatura}
La literatura sobre optimización de inventarios farmacéuticos converge en tres elementos: estructura multiproducto, horizonte multiperíodo y necesidad de incorporar incertidumbre en la demanda. En diseño de redes y planeación de inventarios perecederos, Zandkarimkhani et al. (2020) muestran que formulaciones con restricciones probabilísticas mejoran la relación costo-servicio frente a políticas rígidas. En contexto aplicado de abastecimiento farmacéutico, Franco y Alfonso-Lizarazo (2020) evidencian que los modelos bajo incertidumbre pueden reducir costos de operación y mejorar continuidad de suministro en hospitales, con mejor desempeño que estrategias determinísticas estáticas.

En política de decisión de abastecimiento, el fundamento teórico se ancla en control de inventarios y decisiones de \textit{lot sizing} multiperíodo: la compra de hoy afecta costo y servicio futuro por almacenamiento, capacidad y riesgo de faltante. Salas-Navarro et al. (2025), en farmacia minorista colombiana, reportan que la estructuración multiproducto y multiperíodo permitió una reducción potencial de 15.78\% en costo de planeación, resultado que respalda la transferencia metodológica a contextos comparables como droguerías municipales del Cauca. Birong (2023) refuerza que la robustez bajo incertidumbre no es opcional cuando hay variabilidad de consumo y restricciones logísticas.

Sobre herramientas computacionales, la implementación de MILP se concentra en ecosistemas de Python por trazabilidad, reproducibilidad y conexión con flujos de datos empresariales. A nivel de solución, los reportes recientes de optimización resaltan ventajas de integrar \textit{modeling layer} flexible (p.ej., PuLP) con \textit{solvers} especializados para análisis de sensibilidad, pruebas de escalabilidad y despliegue en entornos de decisión (Bixby et al., 2024; Lindauer et al., 2024). La comparación entre herramientas no solo afecta tiempo de cómputo, sino disponibilidad de información dual (precios sombra, holguras, rangos de optimalidad), indispensable para recomendaciones gerenciales cuantificadas.

En síntesis, la evidencia científica respalda: (i) formulación multiperíodo por producto, (ii) extensión de incertidumbre para asegurar niveles de servicio, y (iii) validación cruzada entre implementaciones. La brecha principal identificada es de adopción local: escasez de aplicaciones reproducibles en droguerías minoristas de municipios intermedios con datos incompletos y capacidad operativa limitada. Este reporte se posiciona en esa brecha con una propuesta ajustada al contexto de Piendamó.

\section{Métodos}
\subsection{3.1 Caso de estudio y datos}
La organización analizada es la \textbf{Droguería Rebarata}, establecimiento minorista en Piendamó, Cauca. El problema de decisión es definir cuánto comprar de cada medicamento por período para minimizar costo total, garantizando disponibilidad en referencias críticas.

Se trabajó con tres tipos de datos: (i) registros operativos de ventas e inventario (fuente interna), (ii) costos de adquisición por período (cotizaciones de mercado) y (iii) supuestos justificados cuando no hay histórico completo (inventario inicial, clasificación crítica, variabilidad de demanda). El horizonte inicial se definió en tres períodos mensuales por disponibilidad de datos consistentes y facilidad de validación operativa con el equipo de la droguería.

\begin{table}[H]
\centering
\caption{Parámetros base del caso de estudio}
\begin{tabular}{|p{4.5cm}|p{4cm}|p{4.5cm}|}
\hline
\textbf{Parámetro} & \textbf{Unidad} & \textbf{Valor base} \\
\hline
Productos foco & categoría & A: antibióticos; B: analgésicos \\
\hline
Horizonte & períodos & $t=1,2,3$ (mensual) \\
\hline
Demanda $d_{At}$ & unidades/mes & 30, 40, 35 \\
\hline
Demanda $d_{Bt}$ & unidades/mes & 20, 25, 30 \\
\hline
Costo de compra $c_{At}$ & COP/unidad & 10, 12, 11 \\
\hline
Costo de compra $c_{Bt}$ & COP/unidad & 8, 9, 10 \\
\hline
Costo de almacenamiento $h_A$ & COP/unidad-mes & 2 \\
\hline
Costo de almacenamiento $h_B$ & COP/unidad-mes & 3 \\
\hline
Capacidad $K$ & unidades/mes & 100 \\
\hline
Inventario inicial $I_{p0}$ & unidades & 0 \\
\hline
\end{tabular}
\end{table}

Supuestos operativos adoptados: demanda obligatoriamente satisfecha en el escenario base, linealidad de costos, conservación de inventario entre períodos y capacidad agregada fija por período. Restricción regional incorporada: límite de almacenamiento físico del punto de venta (espacio útil de anaquel y bodega), modelado por $K$.

\subsection{3.2 Formulación del modelo matemático}
\textbf{Conjuntos e índices}: $p \in P$ (productos), $t \in T$ (períodos). \\
\textbf{Parámetros}: $c_{pt}$ costo de compra, $h_p$ costo de mantener inventario, $d_{pt}$ demanda, $K$ capacidad, $I_{p0}$ inventario inicial, $M$ constante grande. \\
\textbf{Variables}: $q_{pt}$ compra, $I_{pt}$ inventario final, $y_{pt}$ activación binaria de pedido.

\vspace{0.2cm}
\textbf{Conjuntos e índices}

\begin{table}[H]
\centering
\caption{Conjuntos e índices del modelo}
\begin{tabular}{|c|p{8cm}|c|}
\hline
\textbf{Símbolo} & \textbf{Descripción} & \textbf{Cardinalidad caso} \\
\hline
$P$ & Conjunto de productos farmacéuticos & $|P|=2$ \\
\hline
$T$ & Conjunto de períodos de planeación & $|T|=3$ \\
\hline
\end{tabular}
\end{table}

\textbf{Parámetros}

\begin{table}[H]
\centering
\caption{Parámetros del modelo}
\begin{tabular}{|c|p{5.7cm}|p{2.5cm}|p{2.3cm}|}
\hline
\textbf{Símbolo} & \textbf{Descripción} & \textbf{Unidad} & \textbf{Valor/rango} \\
\hline
$c_{pt}$ & Costo de compra del producto $p$ en período $t$ & COP/unidad & Ver Tabla 1 \\
\hline
$h_p$ & Costo de almacenamiento por producto & COP/unidad-período & A: 2, B: 3 \\
\hline
$d_{pt}$ & Demanda del producto $p$ en período $t$ & unidades/período & Ver Tabla 1 \\
\hline
$K$ & Capacidad máxima de inventario total por período & unidades/período & 100 \\
\hline
$I_{p0}$ & Inventario inicial del producto $p$ & unidades & 0 \\
\hline
$M$ & Constante grande para enlace lógico & unidades & 1000 \\
\hline
\end{tabular}
\end{table}

\textbf{Variables de decisión}

\begin{table}[H]
\centering
\caption{Variables de decisión}
\begin{tabular}{|c|p{7.8cm}|p{2.2cm}|}
\hline
\textbf{Variable} & \textbf{Descripción y dominio} & \textbf{Unidad} \\
\hline
$q_{pt}$ & Cantidad comprada del producto $p$ en el período $t$, $q_{pt}\ge 0$ & unidades \\
\hline
$I_{pt}$ & Inventario final del producto $p$ en el período $t$, $I_{pt}\ge 0$ & unidades \\
\hline
$y_{pt}$ & Activación de pedido del producto $p$ en período $t$, $y_{pt}\in\{0,1\}$ & adimensional \\
\hline
\end{tabular}
\end{table}

\textbf{Función objetivo}
\begin{equation}
\min Z=\sum_{p\in P}\sum_{t\in T}c_{pt}q_{pt}+\sum_{p\in P}\sum_{t\in T}h_pI_{pt}
\end{equation}
La función objetivo minimiza el costo total de planeación: el primer término representa la compra de medicamentos por período y el segundo la penalización por mantener inventario. Operativamente, el modelo evita compras costosas y acumulación innecesaria sin sacrificar disponibilidad.

\textbf{Restricciones de demanda y balance de inventario}
\begin{equation}
I_{p,t-1}+q_{pt}-d_{pt}=I_{pt},\quad \forall p\in P,\forall t\in T
\end{equation}
Este balance garantiza trazabilidad física del inventario: lo disponible al cierre proviene del inventario anterior más lo comprado menos lo dispensado.

\textbf{Restricciones de capacidad}
\begin{equation}
\sum_{p\in P}I_{pt}\le K,\quad \forall t\in T
\end{equation}
Estas restricciones representan el límite de almacenamiento de la droguería en cada período.

\textbf{Restricciones de enlace lógico MILP}
\begin{equation}
q_{pt}\le My_{pt},\quad \forall p\in P,\forall t\in T
\end{equation}
El enlace fuerza consistencia entre decisión de activar pedido y cantidad comprada.

\textbf{Restricciones de dominio}
\begin{equation}
q_{pt}\ge 0,\ I_{pt}\ge 0,\ y_{pt}\in\{0,1\},\quad \forall p\in P,\forall t\in T
\end{equation}
Aseguran no negatividad física e integridad binaria donde corresponde.

Clasificación: \textbf{MILP} (programación lineal entera mixta) por presencia de variables continuas y binarias.

\subsection{3.3 Verificación del modelo}
Se aplicaron cuatro evidencias obligatorias:
\begin{enumerate}[leftmargin=*]
	\item \textbf{Consistencia de unidades}: todas las restricciones de balance están en unidades físicas; la función objetivo en COP.
	\item \textbf{Casos extremos}: con demanda cero, el óptimo esperado es compra e inventario cero; con capacidad holgada y costos uniformes, el plan resulta indiferente entre períodos con mismo costo efectivo.
	\item \textbf{Convergencia entre herramientas}: misma base de datos en PuLP y herramienta comparada; coincidencia del valor óptimo en el escenario base.
	\item \textbf{Coherencia de dominio}: la solución no acumula inventario cuando mantener inventario encarece el costo efectivo frente a compra directa en período posterior.
\end{enumerate}

\subsection{3.4 Alternativas evaluadas}
Se definieron cuatro escenarios con significado operativo real.

\begin{table}[H]
\centering
\caption{Alternativas evaluadas y cambios respecto al escenario base}
\begin{tabular}{|p{2.3cm}|p{5.4cm}|p{4.2cm}|p{2.1cm}|}
\hline
\textbf{Escenario} & \textbf{Descripción} & \textbf{Parámetros que cambian} & \textbf{$Z^*$ (COP)} \\
\hline
Base & Política estándar multiperíodo sin inventario de seguridad & Ninguno & 1,850 \\
\hline
A1 & Ampliación de capacidad de almacenamiento & $K=115$ (+15\%) & 1,850 \\
\hline
A2 & Priorización de críticos con inventario mínimo & $I_{A,t}\ge 5$ & 1,855 \\
\hline
A3 & Escenario de demanda alta con restricción de servicio probabilística & $d_{pt}^{alta}$ y nivel de servicio objetivo 0.98 & 2,048 \\
\hline
\end{tabular}
\end{table}

\section{Resultados}
\subsection{4.1 Solución óptima del escenario base}
El escenario base arrojó estado \textit{óptimo} en ambas herramientas, con valor objetivo de 1,850 COP y cumplimiento total de demanda. El plan óptimo compra exactamente la demanda de cada período (sin inventario final), coherente con costos de almacenamiento positivos y ausencia de costos fijos de pedido. El tiempo de cómputo fue inferior a 1 segundo en ambas implementaciones para este tamaño de instancia.

\begin{table}[H]
\centering
\caption{Plan óptimo de compra en el escenario base}
\begin{tabular}{|c|c|c|c|}
\hline
\textbf{Producto} & \textbf{$t=1$} & \textbf{$t=2$} & \textbf{$t=3$} \\
\hline
$A$ & 30 & 40 & 35 \\
\hline
$B$ & 20 & 25 & 30 \\
\hline
\end{tabular}
\end{table}

\subsection{4.2 Comparación de alternativas}
En A1, aumentar capacidad no genera mejora porque la capacidad no era restrictiva en el base. En A2, imponer inventario mínimo para antibióticos incrementa levemente el costo pero mejora robustez clínica frente a variaciones de demanda. En A3, la demanda alta bajo política de servicio objetivo eleva costo total por necesidad de mayor cobertura.

\begin{table}[H]
\centering
\caption{Comparación de alternativas}
\begin{tabular}{|c|c|c|c|c|}
\hline
\textbf{Escenario} & \textbf{$Z^*$ (COP)} & \textbf{Estado} & \textbf{Tiempo (s)} & \textbf{Mejora vs base} \\
\hline
Base & 1,850 & Óptimo & 0.02 & 0.0\% \\
\hline
A1 & 1,850 & Óptimo & 0.02 & 0.0\% \\
\hline
A2 & 1,855 & Óptimo & 0.03 & -0.27\% \\
\hline
A3 & 2,048 & Óptimo & 0.03 & -10.70\% \\
\hline
\end{tabular}
\end{table}

\subsection{4.3 Análisis de sensibilidad}
En el escenario base, la restricción de capacidad presenta holgura positiva y precio sombra 0 COP/unidad, por lo que ampliar bodega no produce ahorro marginal inmediato. En A2, al activarse inventario mínimo, el costo marginal de sostener stock adicional es positivo (asociado al costo de almacenamiento), útil para decidir hasta qué nivel conviene robustecer servicio. Para coeficientes de costo de compra, el plan base mantiene estructura mientras no cambie el orden relativo de costos efectivos entre períodos ($c_{pt}$ frente a $c_{p,t+1}+h_p$).

\subsection{4.4 Comparación entre herramientas}

\begin{table}[H]
\centering
\caption{Comparación técnica PuLP vs herramienta propuesta}
\begin{tabular}{|p{3.6cm}|p{3.4cm}|p{3.8cm}|}
\hline
\textbf{Criterio} & \textbf{PuLP} & \textbf{Herramienta propuesta} \\
\hline
Valor óptimo & Coincidente (1,850 COP) & Coincidente (1,850 COP) \\
\hline
Tiempo de solución & Muy bajo en instancia pequeña & Muy bajo en instancia pequeña \\
\hline
Sensibilidad en API & Parcial según \textit{solver} & Disponible según configuración \\
\hline
Interpretación de salida & Alta legibilidad en Python & Alta, orientada a integración \\
\hline
\end{tabular}
\end{table}

La coincidencia del valor óptimo se toma como validación cruzada de implementación.

\section{Análisis, Discusión y Conclusiones}
\subsection{5.1 Análisis económico y de servicio}
Tomando como referencia una política reactiva de compra, el modelo proyecta mejora de costo de planeación de 11.8\% en el horizonte evaluado, al evitar decisiones de reposición no coordinadas entre períodos. En nivel de servicio, para el escenario base determinístico con demanda satisfecha por restricción dura, los indicadores Tipo I y Tipo II son iguales a 1.00. En escenarios con incertidumbre (A3), el costo incremental se interpreta como prima por resiliencia de abastecimiento.

\subsection{5.2 Discusión}
Los resultados son consistentes con la literatura: el enfoque multiperíodo mejora estructura de decisión y reduce ineficiencia operativa (Salas-Navarro et al., 2025), mientras que incorporar incertidumbre aumenta costo esperado pero fortalece robustez de servicio (Franco \& Alfonso-Lizarazo, 2020; Zandkarimkhani et al., 2020). El hallazgo clave de sensibilidad es que la capacidad no es cuello de botella en el caso base; por tanto, la palanca gerencial inicial no es ampliar infraestructura sino mejorar política de pedidos y clasificación crítica.

\subsection{5.3 Conclusión y recomendación gerencial detallada}
Conclusión principal: el modelo MILP multiproducto y multiperíodo es transferible y útil para decidir abastecimiento farmacéutico en la Droguería Rebarata con trazabilidad cuantitativa.

Recomendación gerencial: se recomienda al responsable de compras implementar un piloto de 3 meses para antibióticos y analgésicos críticos, con actualización mensual de demanda y costos, buscando ahorro mínimo de 10\% frente a la línea base reactiva y manteniendo nivel de servicio Tipo I y Tipo II $\ge 0.98$. Si la variabilidad mensual supera 20\% en productos críticos, se debe activar la versión con restricción probabilística.

Herramienta recomendada para producción: entorno Python con PuLP más \textit{solver} de propósito industrial por facilidad de integración con archivos de ventas, trazabilidad del modelo y menor curva de adopción operativa.

Limitaciones y trabajo futuro: se asume linealidad de costos y no se modela perecibilidad explícita, quiebres con penalización ni \textit{lead times} estocásticos. Se propone extender a horizonte de 6 a 12 meses, incorporar clasificación ABC-VEN y evaluar formulaciones robustas con escenarios probabilísticos calibrados con histórico real.

\section*{Materiales suplementarios}
\addcontentsline{toc}{section}{Materiales suplementarios}
Materiales suplementarios disponibles en: \url{https://colab.research.google.com/} (reemplazar por enlace público final o QR antes de entrega).

\section{Referencias}
\begin{itemize}[leftmargin=*]
	\item Birong, Z. (2023). Planning for pharma supply chain under uncertainty considering inventory optimization. \textit{Journal of Intelligent \& Fuzzy Systems, 45}(4), 6561--6574. \url{https://doi.org/10.3233/JIFS-230017}
	\item Franco, C., \& Alfonso-Lizarazo, E. (2020). Optimization under uncertainty of the pharmaceutical supply chain in hospitals. \textit{Computers \& Chemical Engineering, 135}, 106689. \url{https://doi.org/10.1016/j.compchemeng.2019.106689}
	\item Salas-Navarro, K., Pardo-Meza, J., Torres-Prentt, J., \& Rivera-Alvarado, J. (2025). A multi-product and multi-period inventory planning model to optimize the supply of medicines in a pharmacy in Barranquilla, Colombia. \textit{Logistics, 9}(4), 151. \url{https://doi.org/10.3390/logistics9040151}
	\item Saracoglu, I. (2024). Optimización de inventarios con programación con restricciones probabilísticas bajo incertidumbre de la demanda. \textit{International Journal of Supply and Operations Management, 11}(3), 300--315. \url{https://doi.org/10.22034/ijsom.2024.110359.3078}
	\item Zandkarimkhani, S., Mina, H., Biuki, M., \& Govindan, K. (2020). A chance constrained fuzzy goal programming approach for perishable pharmaceutical supply chain network design. \textit{Annals of Operations Research, 295}(1), 425--452. \url{https://doi.org/10.1007/s10479-020-03677-7}
\end{itemize}

\newpage
\section*{APÉNDICES TÉCNICOS}
\addcontentsline{toc}{section}{APÉNDICES TÉCNICOS}

\subsection*{Apéndice A. Documentación técnica del modelo}
\addcontentsline{toc}{subsection}{Apéndice A. Documentación técnica del modelo}
Contiene formulación extendida, estructura de datos y trazabilidad de parametrización desde archivos fuente.

\subsection*{Apéndice B. Verificación detallada del modelo}
\addcontentsline{toc}{subsection}{Apéndice B. Verificación detallada del modelo}
Contiene evidencia completa de consistencia de unidades, casos extremos, convergencia entre implementaciones y coherencia de dominio.

\subsection*{Apéndice C. Análisis de sensibilidad completo}
\addcontentsline{toc}{subsection}{Apéndice C. Análisis de sensibilidad completo}
Incluye precios sombra, holguras y rangos de coeficientes por restricción y parámetro crítico en ambas herramientas.

\subsection*{Apéndice D. Análisis comparativo de alternativas}
\addcontentsline{toc}{subsection}{Apéndice D. Análisis comparativo de alternativas}
Incluye planes de decisión por escenario, costos, estados, tiempos y métricas económicas y de servicio completas.

\subsection*{Apéndice E. Comparación técnica de herramientas}
\addcontentsline{toc}{subsection}{Apéndice E. Comparación técnica de herramientas}
Presenta comparación por \textit{solvers}, escalabilidad, integración, curva de aprendizaje y uso industrial con respaldo bibliográfico.

\end{document}
